\renewcommand{\bibname}{Список использованных источников}
\addcontentsline{toc}{chapter}{Список использованных источников}
\begingroup
\Urlmuskip=0mu plus 1mu
\urlstyle{same}
\raggedright
\begin{thebibliography}{99}
\bibitem{zaramenskikh} Зараменских Е.\,П. Проектирование информационных систем : учебник и практикум для вузов. - 2-е изд. - Москва : Юрайт, 2025. - 119 с. - ISBN 978-5-534-21418-5. - URL: \url{https://urait.ru/bcode/571331} (дата обращения: 04.11.2026).
\bibitem{chistov} Чистов Д.\,В., Мельников П.\,П., Золотарюк А.\,В., Ничепорук Н.\,Б. Проектирование информационных систем : учебник и практикум для вузов. - 2-е изд., перераб. и доп. - Москва : Юрайт, 2025. - 273 с. - ISBN 978-5-534-20361-5. - URL: \url{https://urait.ru/bcode/560485} (дата обращения: 04.11.2026).
\bibitem{polupanov} Полупанов Д.\,В., Абдюшева С.\,Р., Ефимов А.\,М. Программирование в Python 3 : учебное пособие. - Уфа : БашГУ, 2020. - 164 с. - ISBN 978-5-7477-5230-6. - URL: \url{https://e.lanbook.com/book/179915} (дата обращения: 04.11.2026).
\bibitem{yantsev} Янцев В.\,В. Web-программирование на Python : учебное пособие для вузов. - Санкт-Петербург : Лань, 2022. - 180 с. - ISBN 978-5-8114-9461-3. - URL: \url{https://e.lanbook.com/book/233264} (дата обращения: 04.11.2026).
\bibitem{prokhorenok} Прохоренок Н.\,А., Дронов В.\,А. Python 3 и PyQt 6. Разработка приложений. - Санкт-Петербург : БХВ-Петербург, 2023. - 832 с. - ISBN 978-5-9775-1706-5. - URL: \url{https://library.cbr.ru/catalog/lib/books/404707/} (дата обращения: 04.11.2026).
\bibitem{mdnhttp} MDN Web Docs. Обзор протокола HTTP [Электронный ресурс]. - URL: \url{https://developer.mozilla.org/ru/docs/Web/HTTP/Guides/Overview} (дата обращения: 04.11.2026).
\bibitem{pythonsqlite} Python Software Foundation. sqlite3 - DB-API 2.0 interface for SQLite databases [Электронный ресурс]. - URL: \url{https://docs.python.org/3/library/sqlite3.html} (дата обращения: 04.11.2026).
\bibitem{sqlitecreate} SQLite. CREATE TABLE [Электронный ресурс]. - URL: \url{https://www.sqlite.org/lang_createtable.html} (дата обращения: 04.11.2026).
\bibitem{fastapifirst} FastAPI. First Steps [Электронный ресурс]. - URL: \url{https://fastapi.tiangolo.com/tutorial/first-steps/} (дата обращения: 04.11.2026).
\bibitem{fastapibody} FastAPI. Request Body [Электронный ресурс]. - URL: \url{https://fastapi.tiangolo.com/tutorial/body/} (дата обращения: 04.11.2026).
\bibitem{fastapitesting} FastAPI. Testing [Электронный ресурс]. - URL: \url{https://fastapi.tiangolo.com/tutorial/testing/} (дата обращения: 04.11.2026).
\bibitem{qtdocs} Qt Company. Qt for Python - Documentation [Электронный ресурс]. - URL: \url{https://doc.qt.io/qtforpython-6/} (дата обращения: 04.11.2026).
\bibitem{qttabs} Qt Company. PySide6.QtWidgets.QTabWidget [Электронный ресурс]. - URL: \url{https://doc.qt.io/qtforpython-6/PySide6/QtWidgets/QTabWidget.html} (дата обращения: 04.11.2026).
\bibitem{qtlist} Qt Company. PySide6.QtWidgets.QListWidget [Электронный ресурс]. - URL: \url{https://doc.qt.io/qtforpython-6/PySide6/QtWidgets/QListWidget.html} (дата обращения: 04.11.2026).
\bibitem{qtimer} Qt Company. PySide6.QtCore.QTimer [Электронный ресурс]. - URL: \url{https://doc.qt.io/qtforpython-6/PySide6/QtCore/QTimer.html} (дата обращения: 04.11.2026).
\bibitem{requestsquick} Requests. Quickstart [Электронный ресурс]. - URL: \url{https://requests.readthedocs.io/en/latest/user/quickstart/} (дата обращения: 04.11.2026).
\bibitem{rfchttp} Fielding R., Nottingham M., Reschke J. HTTP Semantics. RFC 9110. - IETF, 2022. - URL: \url{https://datatracker.ietf.org/doc/html/rfc9110} (дата обращения: 04.11.2026).
\bibitem{rfcjson} Bray T. The JavaScript Object Notation (JSON) Data Interchange Format. RFC 8259. - IETF, 2017. - URL: \url{https://datatracker.ietf.org/doc/html/rfc8259} (дата обращения: 04.11.2026).
\bibitem{owaspvalidation} OWASP Cheat Sheet Series. Input Validation Cheat Sheet [Электронный ресурс]. - URL: \url{https://cheatsheetseries.owasp.org/cheatsheets/Input_Validation_Cheat_Sheet.html} (дата обращения: 04.11.2026).
\bibitem{gostflow} ГОСТ 19.701-90 (ИСО 5807-85). Единая система программной документации. Схемы алгоритмов, программ, данных и систем. Обозначения условные и правила выполнения. - Росстандарт. - URL: \url{https://protect.gost.ru/gost/details/0c85e066-f89e-433d-a6d3-b7f746bd0c7c} (дата обращения: 04.11.2026).
\end{thebibliography}
\endgroup
